\documentclass[tikz,border=4pt]{standalone}
\usepackage{ctex}
\usepackage{amsmath,amssymb}
\usetikzlibrary{calc,arrows.meta,decorations.markings}
\begin{document}
\begin{tikzpicture}[
  x={(1cm,0cm)}, y={(0cm,1cm)}, z={(-0.45cm,-0.38cm)},
  line cap=round, line join=round,
  >=Stealth,
  font=\small
]
% 坐标轴
\draw[->] (-0.5,0,0) -- (4,0,0) node[right] {$x$};
\draw[->] (0,-0.5,0) -- (0,3.8,0) node[above] {$y$};
\draw[->] (0,0,-0.3) -- (0,0,3.5) node[above left] {$z$};

% 曲面 Σ（碗形开放曲面 z = 0.3 + 0.2*((x-1.5)^2 + (y-1.5)^2)），用网格
\foreach \ui in {0.3,0.7,...,2.7} {
  \draw[blue!25, thin]
    plot[domain=0.3:2.7, samples=10, variable=\vi]
      ({\ui},{\vi},{0.3+0.15*((\ui-1.5)*(\ui-1.5)+(\vi-1.5)*(\vi-1.5))});
}
\foreach \vi in {0.3,0.7,...,2.7} {
  \draw[blue!25, thin]
    plot[domain=0.3:2.7, samples=10, variable=\ui]
      ({\ui},{\vi},{0.3+0.15*((\ui-1.5)*(\ui-1.5)+(\vi-1.5)*(\vi-1.5))});
}
\node[blue!60!black] at (2.5, 2.5, 1.6) {$\Sigma$};

% 边界 ∂Σ（曲面边界 ，矩形边界提升到曲面 z）— 逆时针
% 四条边
\draw[->, very thick, red!70!black, decoration={markings,
  mark=at position 0.6 with {\arrow{>}}}, postaction={decorate}]
  plot[domain=0.3:2.7, samples=20, variable=\t]
    ({\t},{0.3},{0.3+0.15*((\t-1.5)*(\t-1.5)+(0.3-1.5)*(0.3-1.5))});
\draw[->, very thick, red!70!black, decoration={markings,
  mark=at position 0.6 with {\arrow{>}}}, postaction={decorate}]
  plot[domain=0.3:2.7, samples=20, variable=\t]
    ({2.7},{\t},{0.3+0.15*((2.7-1.5)*(2.7-1.5)+(\t-1.5)*(\t-1.5))});
\draw[->, very thick, red!70!black, decoration={markings,
  mark=at position 0.6 with {\arrow{>}}}, postaction={decorate}]
  plot[domain=2.7:0.3, samples=20, variable=\t]
    ({\t},{2.7},{0.3+0.15*((\t-1.5)*(\t-1.5)+(2.7-1.5)*(2.7-1.5))});
\draw[->, very thick, red!70!black, decoration={markings,
  mark=at position 0.6 with {\arrow{>}}}, postaction={decorate}]
  plot[domain=2.7:0.3, samples=20, variable=\t]
    ({0.3},{\t},{0.3+0.15*((0.3-1.5)*(0.3-1.5)+(\t-1.5)*(\t-1.5))});
\node[red!70!black] at (3.3, 0.5, 0.5) {$\partial\Sigma$};

% 法向量 n（在曲面中央）
\def\u{1.5}\def\v{1.5}
\pgfmathsetmacro{\pz}{0.3+0.15*((\u-1.5)*(\u-1.5)+(\v-1.5)*(\v-1.5))}
\draw[->, very thick, orange] (\u,\v,\pz) -- ++(0,0,1.0)
  node[above, font=\small] {$\mathbf{n}$};

% 右手法则示意：右手图标（简化）
\node[draw=green!50!black, fill=green!10, rounded corners=2pt,
      font=\small, align=left] at (-1, 3, 2.5)
  {\textbf{右手法则}\\
   四指沿 $\partial\Sigma$\\
   拇指指向 $\mathbf{n}$};

% 曲面上若干旋度向量 ∇×F
\foreach \xx/\yy in {1.0/1.0, 2.0/1.0, 1.0/2.0, 2.0/2.0} {
  \pgfmathsetmacro{\zz}{0.3+0.15*((\xx-1.5)*(\xx-1.5)+(\yy-1.5)*(\yy-1.5))}
  \draw[->, purple!70, thick] (\xx,\yy,\zz) -- ++(0,0,0.5);
}
\node[purple!70!black, font=\small] at (3.5, 1.5, 1.3) {$\nabla\times\mathbf{F}$（旋度场）};

% 公式
\node[font=\small, draw=gray!50, fill=gray!10, rounded corners=2pt, align=left]
  at (3.3, 3.5, -0.5)
  {\textbf{Stokes 旋度定理}\\[2pt]
   $\displaystyle\iint_\Sigma(\nabla\!\times\!\mathbf{F})\cdot d\mathbf{S}=\!\!\oint_{\partial\Sigma}\!\mathbf{F}\cdot d\mathbf{r}$\\[2pt]
   "曲面旋度积分 = 边界环量"};
\end{tikzpicture}
\end{document}
